% Jednotlivé vzory jsou zapsány jako seznamy černých pixelů x/y.
% Mřížka má 10 sloupců a 12 řádků, tedy 120 neuronů.

\def\patternThree{
    2/11,3/11,4/11,5/11,6/11,7/11,
    1/10,2/10,7/10,8/10,
    7/9,8/9,
    6/8,7/8,8/8,
    4/7,5/7,6/7,7/7,
    4/6,5/6,6/6,7/6,8/6,
    7/5,8/5,
    7/4,8/4,
    7/3,8/3,
    1/2,2/2,7/2,8/2,
    2/1,3/1,4/1,5/1,6/1,7/1
}

\def\patternNoisy{
    0/11,2/11,3/11,4/11,5/11,7/11,
    7/10,8/10,
    5/9,8/9,
    0/8,2/8,5/8,6/8,7/8,8/8,
    0/7,4/7,5/7,6/7,7/7,9/7,
    0/6,6/6,7/6,
    0/5,8/5,9/5,
    1/4,7/4,9/4,
    5/3,
    1/2,7/2,8/2,9/2,
    4/1,5/1,7/1,
    3/0
}

\def\patternIntermediateA{
    2/11,3/11,4/11,5/11,6/11,7/11,
    1/10,2/10,7/10,8/10,
    7/9,8/9,
    6/8,7/8,8/8,
    1/7,4/7,5/7,7/7,
    4/6,5/6,6/6,8/6,
    7/5,9/5,
    1/4,7/4,8/4,
    7/3,8/3,
    1/2,2/2,7/2,8/2,
    2/1,3/1,5/1,6/1,7/1,9/1
}

\def\patternIntermediateB{
    2/11,3/11,4/11,5/11,6/11,7/11,
    1/10,2/10,7/10,8/10,
    7/9,8/9,
    7/8,8/8,
    4/7,5/7,6/7,7/7,
    4/6,5/6,6/6,7/6,8/6,
    7/5,8/5,
    7/4,8/4,
    7/3,8/3,9/3,
    1/2,2/2,7/2,8/2,
    2/1,3/1,4/1,5/1,6/1,7/1
}

% #1, #2 ... posunutí obrázku
% #3     ... seznam černých pixelů
% #4     ... nadpis
% #5     ... barva nadpisu
\newcommand{\hopfieldpattern}[5]{%
    \begin{scope}[shift={(#1,#2)}]

        \node[
            task,
            fill=#5!15,
            minimum width=23mm,
            minimum height=7mm,
            font=\scriptsize
        ] at (0.75,2.5) {#4};

        \fill[white]
            (0,0) rectangle (1.5,1.8);

        \draw[
            step=0.15cm,
            very thin,
            black!25
        ] (0,0) grid (1.5,1.8);

        % Nejprve se rozbalí makro se seznamem pixelů
        \edef\drawpatternloop{%
            \noexpand\foreach
            \noexpand\x/\noexpand\y
            in {#3}%
        }%

        % Teprve poté se provede cyklus
        \drawpatternloop {%
            \fill[black]
                ({0.15*\x},{0.15*\y})
                rectangle ++(0.15,0.15);
        }%

        \draw[thick]
            (0,0) rectangle (1.5,1.8);

    \end{scope}
}